\section{Algorithm Cheat-Sheet}
\label{sec:algocheat}

A compact one-row-per-algorithm reference for the ``describe / sketch / MCQ''
question type. Detailed treatments are in Sections~7--11; this is the revision
sheet. Sketch boxes follow for the algorithms where a drawing matters.

\renewcommand{\arraystretch}{1.25}
\begin{longtable}{@{}>{\raggedright\arraybackslash}p{2.5cm}
                     >{\raggedright\arraybackslash}p{4.3cm}
                     >{\raggedright\arraybackslash}p{3.6cm}
                     >{\raggedright\arraybackslash}p{4.6cm}@{}}
\toprule
\textbf{Algorithm} & \textbf{Core idea} & \textbf{Key hyperparameters} &
\textbf{Key properties / pitfalls}\\
\midrule
\endfirsthead
\toprule
\textbf{Algorithm} & \textbf{Core idea} & \textbf{Key hyperparameters} &
\textbf{Key properties / pitfalls}\\
\midrule
\endhead
\bottomrule
\endfoot

Recursive subdivision (mosaic) & Split a rectangle by category proportions,
alternating axes per variable & Variable order; spacing & Tile area = joint
frequency; order changes layout; reads marginals/conditionals\\

Recursive subdivision (treemap) & Recursively nest rectangles sized by a
hierarchy's leaf values & Tiling (slice-and-dice, squarified); layout order &
Area = value; squarified gives good aspect ratios; deep hierarchies hard to
read; area is a weak channel\\

Just-noticeable-difference (JND) & Smallest stimulus change a viewer can
detect & (Weber fraction $k$) & Weber's law: $\Delta I / I = k$; motivates
perceptually uniform colour/size scales\\

Scagnostics & Graph-theoretic measures summarizing a scatterplot's shape &
(Binning grid; measure set) & Outlying, skewed, clumpy, sparse, striated,
convex, skinny, stringy, monotonic; used to rank SPLOM panels\\

Class consistency & Quality metric: do class labels agree with spatial
structure of a projection? & Distance/centroid criterion & Higher = classes
form separable clusters; used to evaluate/rank DR layouts\\

PCA & Linear projection onto orthogonal directions of max variance &
\#components $k$ & Linear, deterministic; preserves global variance; axes
interpretable; misses nonlinear structure\\

MDS & Embed so low-D distances match high-D (dis)similarities & metric vs.\
non-metric; \#dims; stress & Preserves pairwise distances; classical MDS
$\approx$ PCA; can be slow; stress measures distortion\\

t-SNE & Match neighbour probability distributions (high-D vs.\ low-D),
minimize KL & \term{perplexity}; learning rate; iterations; seed &
Nonlinear, stochastic; preserves local clusters; \emph{cluster sizes \&
between-cluster distances not meaningful}\\

UMAP & Build a fuzzy $k$-NN graph, optimize a low-D layout of it &
\term{$n\_neighbors$}; \term{min\_dist}; \#dims; seed & Nonlinear,
stochastic; faster than t-SNE; better global structure; sizes/distances still
unreliable\\

Hexagon binning & Aggregate points into hexagonal bins, encode count by
colour & bin size / \#bins & Fixes overplotting (perceptual scalability);
hexagons tile uniformly, less axis bias than squares\\

Kernel density estimation & Sum smooth kernels over points to estimate a
density & \term{bandwidth}; kernel shape & Bandwidth too small = spiky, too
large = oversmoothed; $O(n)$ per eval -- costly on tall data\\

Multi-dim tables / data cubes & Pre-aggregate measures over dimension
combinations for fast queries & dimensions; hierarchies; measures &
Interactive scalability: slice/dice/roll-up/drill-down in $O(1)$; storage
grows with dimensions\\

Parallel aggregation & Distribute aggregation over partitions, combine
partial results & \#partitions & Scales aggregation to tall data; basis of
data-cube construction; needs associative aggregates\\

Feature vectors (neuron activations) & Represent an input by a layer's
activation vector & layer choice; (DR for display) & Used to inspect what a
network ``sees''; high-D, usually projected with t-SNE/UMAP\\

Local perturbations (LIME) & Perturb one input, fit a simple local surrogate &
\#samples; kernel width; surrogate & Model-agnostic, \emph{local} explanation;
unstable to sampling; not globally faithful\\

Partial dependence (PDP) & Average model output while sweeping one feature
over its range & feature(s); grid resolution & Global, model-agnostic; assumes
feature independence; can hide interactions (ICE curves fix this)\\

Counterfactual explanation & Smallest input change that flips the prediction &
distance metric; target class; sparsity & ``Change X to get outcome Y'';
should be close, sparse, actionable, plausible; many valid counterfactuals\\

Force-directed layout & Springs on edges + repulsion on nodes, minimize energy
& spring constant; repulsion; iterations; temperature & Good for small/medium
graphs; stochastic; $O(n^2)$ per step (Barnes--Hut speeds it); can get stuck
in local minima\\

Layered (Sugiyama) layout & Assign nodes to layers, order to cut crossings,
assign coordinates & layer assignment; crossing-reduction passes & For
DAGs/hierarchies; minimizes edge crossings; long edges via dummy nodes\\

Graph morphing & Smoothly interpolate between two layouts of a graph &
interpolation; duration & Preserves the \term{mental map} during transitions;
animate node positions\\

Foresighted layout & Lay out a dynamic graph anticipating future states for
stability & time window; stability weight & Trades layout quality vs.\
stability across time steps; keeps the mental map over a sequence\\

\end{longtable}
\renewcommand{\arraystretch}{1.0}

% ---------------------------------------------------------------------
\begin{algobox}[PCA -- 1D projection sketch]
Centre the data, find the direction of maximum variance (first eigenvector of
the covariance matrix), and project points onto it.
\begin{center}
\begin{tikzpicture}[scale=1.0,>=Stealth]
  % tilted point cloud (rotated ellipse), PC1 along 30 deg
  \begin{scope}[rotate=30]
    \draw[rulec] (0,0) ellipse [x radius=2.2cm, y radius=0.8cm];
    % PC1 arrow along long axis
    \draw[->,tuwblue,very thick] (-2.6,0) -- (2.6,0) node[above,font=\footnotesize] {PC1};
    % PC2 orthogonal, discarded
    \draw[->,gray!60,thick] (0,-1.0) -- (0,1.2) node[right,font=\footnotesize] {PC2};
  \end{scope}
  % a few sample points (in rotated frame coords, drawn via scope)
  \begin{scope}[rotate=30]
    \foreach \px/\py in {-1.6/0.45, -0.6/-0.3, 0.4/0.35, 1.4/-0.25, 0.9/0.5} {
      \fill[tuwdark] (\px,\py) circle (1.3pt);
      % perpendicular onto PC1 (the x-axis of rotated frame)
      \draw[dashed,pitred,thin] (\px,\py) -- (\px,0);
      \fill[pitred] (\px,0) circle (1pt);
    }
  \end{scope}
\end{tikzpicture}
\end{center}
\sketchnote{Draw an elongated 2D point cloud. Draw an arrow through its long
axis (PC1). Drop a perpendicular from each point onto that line; the feet of
the perpendiculars are the 1D coordinates. PC2 is orthogonal and discarded.}
\end{algobox}

\begin{algobox}[Recursive subdivision -- mosaic / treemap sketch]
Both split a rectangle proportionally and recurse.
\begin{center}
\begin{tikzpicture}[scale=1.0,line join=round]
  % --- Mosaic: width by A, then each column's height by B ---
  \begin{scope}
    \node[above,font=\footnotesize\bfseries,tuwdark] at (1.4,2.0) {Mosaic};
    % three columns of widths 1.4, 0.9, 0.5 (A proportions); total 2.8 wide, 1.8 tall
    % column 1
    \fill[tuwblue!25] (0,0) rectangle (1.4,1.1);
    \fill[tuwblue!12] (0,1.1) rectangle (1.4,1.8);
    % column 2
    \fill[keygreen!25] (1.4,0) rectangle (2.3,0.7);
    \fill[keygreen!12] (1.4,0.7) rectangle (2.3,1.8);
    % column 3
    \fill[examorange!25] (2.3,0) rectangle (2.8,1.3);
    \fill[examorange!12] (2.3,1.3) rectangle (2.8,1.8);
    % borders
    \draw[white,thin] (1.4,0)--(1.4,1.8) (2.3,0)--(2.3,1.8)
                      (0,1.1)--(1.4,1.1) (1.4,0.7)--(2.3,0.7) (2.3,1.3)--(2.8,1.3);
    \draw[tuwdark,thin] (0,0) rectangle (2.8,1.8);
  \end{scope}
  % --- Treemap: nested rectangles ---
  \begin{scope}[xshift=3.6cm]
    \node[above,font=\footnotesize\bfseries,tuwdark] at (1.4,2.0) {Treemap};
    \draw[tuwdark,thin] (0,0) rectangle (2.8,1.8);
    \fill[tuwblue!20] (0.04,0.04) rectangle (1.6,1.76);
    \fill[keygreen!20] (1.64,0.94) rectangle (2.76,1.76);
    \fill[examorange!20] (1.64,0.04) rectangle (2.76,0.9);
    % recurse inside first child
    \draw[tuwdark!60,thin] (0.04,0.04) rectangle (1.6,1.76);
    \draw[tuwdark!60,thin] (0.04,0.9) rectangle (0.85,1.76);
    \draw[tuwdark!60,thin] (0.04,0.04) rectangle (1.6,0.9);
    \draw[tuwdark!60,thin] (1.64,0.94) rectangle (2.76,1.76);
    \draw[tuwdark!60,thin] (1.64,0.04) rectangle (2.76,0.9);
  \end{scope}
\end{tikzpicture}
\end{center}
\sketchnote{Mosaic: split the full width by variable~A's proportions, then
split each resulting column's height by variable~B's conditional proportions;
tile area = joint frequency. Treemap: split the rectangle by the top-level
children's summed values, then recurse inside each child; leaf area = value.}
\end{algobox}

\begin{algobox}[Hexagon binning sketch]
\begin{center}
\begin{tikzpicture}[scale=1.0]
  % hexagon parameters: minimum size gives the diameter (point-to-point)
  \def\hs{0.7cm}     % hex minimum size
  \def\dx{0.606}     % horizontal spacing approx = 0.866*0.7
  \def\dy{0.525}     % vertical row offset = 0.75*0.7
  % grid of hexagons, offset rows; shade two by "count"
  \foreach \row in {0,1,2} {
    \foreach \col in {0,1,2,3} {
      \pgfmathsetmacro{\cx}{\col*0.866*0.7 + mod(\row,2)*0.433*0.7}
      \pgfmathsetmacro{\cy}{\row*0.525}
      % choose a couple of shaded ones
      \ifnum\row=1 \ifnum\col=1 \def\fc{tuwblue!55}\else\def\fc{none}\fi
      \else \ifnum\row=1 \def\fc{none}\else \def\fc{none}\fi \fi
      \node[regular polygon, regular polygon sides=6, draw=tuwblue!50, thin,
            minimum size=\hs, inner sep=0pt] at (\cx,\cy) {};
    }
  }
  % shade specific hexagons by count
  \node[regular polygon, regular polygon sides=6, draw=tuwblue!50, thin,
        fill=tuwblue!55, minimum size=\hs, inner sep=0pt]
        at ({1*0.866*0.7 + 0.433*0.7},{1*0.525}) {};
  \node[regular polygon, regular polygon sides=6, draw=tuwblue!50, thin,
        fill=tuwblue!25, minimum size=\hs, inner sep=0pt]
        at ({2*0.866*0.7},{2*0.525}) {};
  % scatter dots overlaid
  \foreach \px/\py in {0.5/0.5,0.7/0.6,0.9/0.45,1.0/0.55,0.6/0.4,
                       1.05/1.0,1.3/0.95,1.7/1.05,0.4/1.0,2.0/0.4} {
    \fill[tuwdark] (\px,\py) circle (1pt);
  }
\end{tikzpicture}
\end{center}
\sketchnote{Overlay a hexagonal grid on the scatterplot, count points per
hexagon, and fill each hexagon with a sequential colour by count. Replaces a
dense black blob with a readable density field.}
\end{algobox}

\begin{algobox}[KDE -- bandwidth sketch]
\begin{center}
\begin{tikzpicture}[scale=1.0,>=Stealth]
  % axis
  \draw[->] (0,0) -- (6,0) node[right,font=\footnotesize] {$x$};
  \draw[->] (0,0) -- (0,2.2) node[above,font=\footnotesize] {density};
  % data points (ticks)
  \foreach \dx in {1.0,1.4,1.7,3.6,4.0,4.4} {
    \draw[tuwdark,thick] (\dx,0) -- (\dx,-0.12);
  }
  % small individual kernels (bumps) over each point
  \foreach \dx in {1.0,1.4,1.7,3.6,4.0,4.4} {
    \draw[keygreen!70,thin,smooth] plot[domain=\dx-0.7:\dx+0.7,samples=20]
      (\x, {0.35*exp(-((\x-\dx)*(\x-\dx))/(2*0.13))});
  }
  % bolder summed (good bandwidth) bimodal curve above
  \draw[tuwblue,very thick,smooth] plot[domain=0.1:5.6,samples=80]
    (\x, {0.55*exp(-((\x-1.37)*(\x-1.37))/(2*0.22))
          + 0.55*exp(-((\x-4.0)*(\x-4.0))/(2*0.22)) + 0.05});
  \node[tuwblue,font=\footnotesize] at (5.4,1.65) {sum};
\end{tikzpicture}
\end{center}
\sketchnote{Place a small bump (kernel) on each data point on the x-axis and
sum them into a smooth curve. Show three curves: tiny bandwidth (spiky,
overfit), good bandwidth (smooth bimodal), huge bandwidth (one flat hump,
oversmoothed).}
\end{algobox}

\begin{algobox}[Force-directed layout sketch]
\begin{center}
\begin{tikzpicture}[scale=1.0,>=Stealth,
    nd/.style={circle,draw=tuwblue,fill=tuwblue!15,thick,minimum size=6mm,inner sep=0pt,font=\footnotesize}]
  % five nodes
  \node[nd] (a) at (0,1.4) {1};
  \node[nd] (b) at (1.4,2.0) {2};
  \node[nd] (c) at (1.8,0.6) {3};
  \node[nd] (d) at (3.3,1.6) {4};
  \node[nd] (e) at (3.4,0.2) {5};
  % edges (one labelled spring)
  \draw[tuwdark] (a) -- node[above,font=\scriptsize,sloped] {spring} (b);
  \draw[tuwdark] (a) -- (c);
  \draw[tuwdark] (b) -- (c);
  \draw[tuwdark] (b) -- (d);
  \draw[tuwdark] (d) -- (e);
  \draw[tuwdark] (c) -- (e);
  % repulsion arrows (push apart) between two non-adjacent nodes
  \draw[->,pitred,thin] (b) -- ($(b)!0.35!(e)$);
  \draw[->,pitred,thin] (e) -- ($(e)!0.35!(b)$);
  \node[pitred,font=\scriptsize] at ($(b)!0.5!(e)+(0.55,0)$) {repel};
\end{tikzpicture}
\end{center}
\sketchnote{Nodes as charged particles repel each other; edges act as springs
pulling connected nodes together. Iterate until the energy settles. Draw a
small graph with well-separated nodes and short, roughly equal-length edges,
clusters visibly grouped.}
\end{algobox}

\begin{algobox}[Layered (Sugiyama) layout sketch]
\begin{center}
\begin{tikzpicture}[scale=1.0,>=Stealth,
    nd/.style={circle,draw=tuwblue,fill=tuwblue!15,thick,minimum size=5.5mm,inner sep=0pt,font=\scriptsize}]
  % three layer baselines (dashed)
  \foreach \y/\lbl in {2.4/L1, 1.2/L2, 0/L3} {
    \draw[dashed,rulec] (-0.4,\y) -- (5.0,\y);
    \node[gray,font=\scriptsize] at (5.3,\y) {\lbl};
  }
  % layer 1
  \node[nd] (a) at (1.0,2.4) {a};
  \node[nd] (b) at (3.2,2.4) {b};
  % layer 2
  \node[nd] (c) at (0.6,1.2) {c};
  \node[nd] (d) at (2.2,1.2) {d};
  \node[nd] (e) at (3.8,1.2) {e};
  % layer 3
  \node[nd] (f) at (1.4,0) {f};
  \node[nd] (g) at (3.4,0) {g};
  % downward edges, minimal crossings
  \draw[->,tuwdark] (a) -- (c);
  \draw[->,tuwdark] (a) -- (d);
  \draw[->,tuwdark] (b) -- (e);
  \draw[->,tuwdark] (b) -- (d);
  \draw[->,tuwdark] (c) -- (f);
  \draw[->,tuwdark] (d) -- (f);
  \draw[->,tuwdark] (d) -- (g);
  \draw[->,tuwdark] (e) -- (g);
\end{tikzpicture}
\end{center}
\sketchnote{Draw horizontal layers (rows). Place each node in a layer by its
topological rank; route all edges downward; reorder nodes within layers to
minimize crossings; insert dummy nodes for edges spanning multiple layers.}
\end{algobox}

\begin{algobox}[Counterfactual explanation sketch]
\begin{center}
\begin{tikzpicture}[scale=1.0,>=Stealth]
  % feature space axes
  \draw[->] (0,0) -- (4.6,0) node[right,font=\footnotesize] {$x_1$};
  \draw[->] (0,0) -- (0,3.0) node[above,font=\footnotesize] {$x_2$};
  % curved decision boundary
  \draw[tuwblue,very thick,smooth] plot[domain=0.2:4.2,samples=60]
    (\x, {1.4 + 0.8*sin((\x*55)) });
  % region labels
  \node[pitred,font=\scriptsize] at (0.9,0.5) {rejected};
  \node[keygreen,font=\scriptsize] at (3.6,2.6) {accepted};
  % query point (rejected side)
  \fill[pitred] (1.6,0.7) circle (2pt);
  \node[pitred,font=\scriptsize,below] at (1.6,0.6) {query};
  % counterfactual: nearest point across boundary
  \fill[keygreen] (2.05,1.85) circle (2pt);
  \node[keygreen,font=\scriptsize,right] at (2.1,1.95) {counterfactual};
  % arrow from query to counterfactual
  \draw[->,tuwdark,thick] (1.6,0.7) -- (2.05,1.85);
\end{tikzpicture}
\end{center}
\sketchnote{Plot the decision boundary in feature space with the query point on
the ``rejected'' side. The counterfactual is the nearest point across the
boundary; draw the short arrow from query to that point -- the minimal,
sparse change that flips the outcome.}
\end{algobox}

\begin{algobox}[Partial dependence sketch]
\begin{center}
\begin{tikzpicture}[scale=1.0,>=Stealth]
  \draw[->] (0,0) -- (5.0,0) node[right,font=\footnotesize] {feature value};
  \draw[->] (0,0) -- (0,2.6) node[above,font=\footnotesize,align=center] {avg.\\prediction};
  % smooth response curve (rising then plateau)
  \draw[tuwblue,very thick,smooth] plot[domain=0.2:4.6,samples=60]
    (\x, {0.4 + 1.7/(1+exp(-2.2*(\x-2.2)))});
\end{tikzpicture}
\end{center}
\sketchnote{x-axis = one feature's value, y-axis = average model prediction.
For each grid value of the feature, set it for all rows, predict, and average.
Plot the resulting curve; its slope shows the feature's marginal effect.}
\end{algobox}
